% De acuerdo con la NTC 1486, la introducción presenta la importancia, el origen, los objetivos, el alcance, las limitaciones y la metodología del trabajo, sin anticipar resultados ni conclusiones (2 a 4 páginas).

Buscar empleo remoto exige hoy recorrer varias plataformas, cada una con su propio catálogo, sus propios filtros y un mecanismo de búsqueda que, en lo esencial, compara las palabras de la consulta con las palabras de la oferta. Ese mecanismo funciona cuando quien busca y quien publica describen el cargo con los mismos términos, y falla en cuanto no lo hacen. "Backend Developer" y "Software Engineer", "JS" y "JavaScript", "trabajo con datos" y "Analytics Engineer" pueden referirse a lo mismo sin compartir una sola palabra. La literatura de recuperación de información conoce este problema como \emph{desajuste de vocabulario}. La industria del reclutamiento a gran escala lo ha abordado desde hace casi una década con representaciones semánticas densas \citep{HaThuc2017,Shen2024}, pero las herramientas a disposición de quien busca empleo apenas lo reflejan.

La generación aumentada por recuperación (\textit{Retrieval-Augmented Generation}, RAG) \citep{Lewis2020} ofrece una alternativa. En lugar de exigir coincidencias exactas, recupera los fragmentos de información semánticamente más cercanos a la consulta y delega en un modelo de lenguaje la síntesis de una respuesta fundamentada en ellos. Aplicado a la búsqueda de empleo, este enfoque permite formular preguntas que un panel de filtros no puede expresar: qué ofertas de \textit{backend} en el propio \textit{stack} superan cierto salario, o cuáles de las ofertas guardadas se ajustan a la experiencia descrita en un currículo.

Este trabajo parte de JobAI, una plataforma de búsqueda de empleo remoto construida previamente por el autor, que agrega ofertas de tres fuentes públicas y expone un copiloto conversacional basado en un pipeline RAG. JobAI existe y funciona, pero nunca se había evaluado de forma sistemática. No se sabía, con evidencia cuantitativa, en qué medida su recuperación semántica supera a la búsqueda por palabra clave que pretende reemplazar. Tampoco se sabía qué tan fieles son sus respuestas a los datos recuperados, ni si la decisión inicial de almacenar los vectores en PostgreSQL mediante pgvector era adecuada frente a una base de datos vectorial dedicada. Esos tres vacíos dan origen a la investigación.

\section*{Objetivos}

El objetivo general es desarrollar y evaluar una plataforma de búsqueda de empleo remoto que utiliza un pipeline RAG para permitir la búsqueda semántica y conversacional sobre ofertas laborales, superando las limitaciones de precisión de la búsqueda por palabras clave. Para alcanzarlo se plantean cuatro objetivos específicos:
\begin{enumerate}
    \item Caracterizar, a partir de la literatura, las limitaciones de los sistemas actuales de búsqueda de empleo y las oportunidades que ofrece RAG.
    \item Diseñar y documentar la arquitectura del pipeline, incluidos la recuperación con filtros, la búsqueda híbrida, el reordenamiento y la migración del almacenamiento vectorial a una base de datos dedicada.
    \item Definir una metodología de evaluación con métricas cuantitativas.
    \item Evaluar el artefacto frente a un enfoque de búsqueda tradicional por palabras clave.
\end{enumerate}
El Capítulo~\ref{cap:problema} desarrolla el planteamiento completo.

\section*{Metodología}

El trabajo se enmarca en el paradigma de Design Science Research \citep{Hevner2004,Peffers2007}, que concibe la investigación como el diseño y la evaluación rigurosa de un artefacto que resuelve un problema relevante. Como el artefacto ya existía, se adopta el punto de entrada del proceso centrado en el diseño y desarrollo, y el trabajo se plantea como una nueva iteración del ciclo de diseño-evaluación. Esa iteración tiene dos partes. La de diseño consiste en corregir y migrar el almacenamiento vectorial y en revisar el artefacto frente a su documentación. La de evaluación incluye tres mediciones:
\begin{itemize}
    \item La calidad de la recuperación, frente a dos líneas base léxicas, sobre un corpus real congelado y un conjunto de consultas en cinco categorías.
    \item La fidelidad de las respuestas generadas, con las métricas de RAGAS \citep{Es2024}.
    \item La latencia del pipeline y un benchmark de pgvector frente a Qdrant, con datos reales y sintéticos, con y sin filtros estructurados.
\end{itemize}
Los juicios de relevancia los asigna un modelo de lenguaje, validado contra una muestra de juicios humanos.

\section*{Alcance y aporte}

El estudio se limita a JobAI y a su dominio: ofertas de empleo remoto en inglés, mayoritariamente del sector tecnológico, obtenidas de Remotive, Jobicy y Findwork. Evalúa la calidad técnica del pipeline y no la experiencia de usuarios reales, y compara dos representantes de los enfoques de almacenamiento vectorial, una extensión sobre un motor relacional y un motor dedicado, sin pretender una revisión exhaustiva del mercado.

El aporte es doble. En lo aplicado, ofrece una evaluación cuantitativa e integrada de un sistema RAG completo en producción en un dominio donde la literatura de RAG y la de emparejamiento de empleo rara vez se han cruzado \citep{RojasGaleano2022}. En lo metodológico, documenta cómo aplicar Design Science Research a un artefacto preexistente, incluido el hallazgo de que la propia preparación de la evaluación reveló defectos que la habrían invalidado.

\section*{Estructura del documento}

Además de esta introducción, el documento tiene seis capítulos:
\begin{itemize}
    \item El Capítulo~\ref{cap:problema} plantea el problema, los objetivos y el marco conceptual.
    \item El Capítulo~\ref{cap:estado-arte} revisa el estado del arte en seis líneas: RAG, emparejamiento de empleo, límites de la búsqueda léxica, Design Science Research, bases de datos vectoriales y clasificación híbrida de consultas.
    \item El Capítulo~\ref{cap:metodologia} detalla la metodología de evaluación.
    \item El Capítulo~\ref{cap:artefacto} documenta la arquitectura de JobAI y la iteración de diseño realizada.
    \item El Capítulo~\ref{cap:resultados} presenta y discute los resultados.
    \item El último capítulo responde a cada objetivo específico y propone trabajo futuro.
\end{itemize}
